Neural networks have become a promising data-driven tool for modelling physical systems, offering an alternative approach for approximating dynamical behaviour without explicitly solving the governing equations. This work investigates the ability of a multilayer perceptron (MLP) to approximate the temporal evolution of dynamical systems with increasing complexity. Six representative systems were considered: the harmonic oscillator, damped harmonic oscillator, forced harmonic oscillator, Duffing oscillator, simple pendulum, and double pendulum.

A common neural network architecture was trained using time as the only input variable, while the corresponding physical state was used as the output. Parameter sweep experiments were performed for each system to evaluate how variations in physical parameters influence prediction accuracy. Model performance was assessed using the mean squared error (MSE), mean absolute error (MAE), maximum error, and coefficient of determination $(R^2)$.

The results show that the MLP can accurately approximate the investigated systems across a wide range of explored conditions. The experiments indicate that prediction accuracy is influenced more strongly by the temporal complexity of the generated trajectories than by the presence of nonlinear terms in the governing equations alone. Higher oscillation frequencies, longer prediction intervals, and sensitivity to initial conditions were found to increase prediction errors.

These findings demonstrate the capability of simple feedforward neural networks for approximating deterministic dynamical systems and provide insight into the factors that limit their predictive accuracy.